\chapter{Aksjomaty ZF, ZFC i konstrukcje zbiorowe}\label{ch:zf-axioms}
\chaptermark{Aksjomaty ZF i ZFC}

\section{Dlaczego nie każdy opis wyznacza zbiór?}
We wcześniejszych rozdziałach pracowaliśmy ze zbiorami w elementarnym
języku matematyki. Teraz pytamy, co dokładnie wolno w tym języku
\emph{przyjąć} i jakie zbiory możemy \emph{skonstruować} z aksjomatów.
Naiwna zasada „każda własność wyznacza zbiór wszystkich obiektów,
które ją mają” prowadzi do sprzeczności. Gdyby istniał zbiór
$R=\{x:x\notin x\}$, mielibyśmy
$R\in R\iff R\notin R$. Nie jest to problem z prawami logiki;
problemem jest nieograniczone tworzenie zbiorów.

\begin{definicja}[Język teorii mnogości]\label{def:zf-language}
Językiem teorii Zermela--Fraenkla jest \emph{język pierwszego rzędu}
z jednym dwuargumentowym symbolem relacyjnym $\in$ i równością.
Jego zmienne $x,y,z,\ldots$ przebiegają wyłącznie po \emph{zbiorach}.
Zdanie $x\in y$ oznacza, że zbiór $x$ jest elementem zbioru $y$.
Teoria $\mathrm{ZF}$ składa się z opisanych niżej aksjomatów
i schematów aksjomatów. Teoria $\mathrm{ZFC}$ to $\mathrm{ZF}$
wraz z \emph{aksjomatem wyboru}, który wprowadzimy i omówimy
w rozdziale~\ref{ch:axiom-choice}.
\end{definicja}


\section{Jak czytać aksjomaty: znaczenie zmiennych i skrótów}\label{sec:zf-reading-notation}

\emph{Najważniejsza różnica względem zwykłych podręczników z analizy:}
w formalnym języku ZF nie ma osobnej kategorii „elementów”
i osobnej kategorii „zbiorów”. \emph{Wszystkie zmienne przebiegają po zbiorach}.
Litery $x,a,b,p,u$ nie oznaczają innych rodzajów obiektów
niż wielkie litery $A,B,X$. Zmienne pisane małą literą
są tylko \emph{konwencją redakcyjną}. Dlatego
\[
             x\in a
\]
znaczy dokładnie: „zbiór oznaczony przez $x$ jest elementem
zbioru oznaczonego przez $a$”. Nie oznacza to, że
$x$ jest małym obiektem, a $a$ małym zbiorem.
Zamiast $x\in a$ moglibyśmy napisać $x\in A$,
$u\in B$ albo $A\in B$, gdyby symbole były odpowiednio
zdefiniowane i związane kwantyfikatorami. Wielkość liter
nie wpływa na prawdziwość zdania.

\begin{przyklad}[Element sam jest zbiorem]
Niech $a=\{\varnothing,\{\varnothing\}\}$
i $x=\{\varnothing\}$. Wtedy $x\in a$, ponieważ
$x$ jest jednym z dwóch elementów zbioru $a$.
Jednocześnie $x$ \emph{też jest zbiorem}, a jego
jedynym elementem jest $\varnothing$.
Mamy więc $\varnothing\in x$ oraz $x\in a$,
również $x\subseteq a$, ponieważ $\varnothing\in a$; jednak
te dwa fakty wynikają z osobnych sprawdzeń. W ogólności
$x\in a$ nie implikuje $x\subseteq a$.
Dla kontrprzykładu do implikacji $x\in a\Rightarrow x\subseteq a$
weźmy $a=\{\{\{\varnothing\}\}\}$ i $x=\{\{\varnothing\}\}$.
Wtedy $x\in a$, ale $\{\varnothing\}\in x$ oraz
$\{\varnothing\}\notin a$, więc $x\not\subseteq a$.
\end{przyklad}

\begin{definicja}[Słownik notacji aksjomatów]\label{def:zf-notation-conventions}
W całym rozdziale używamy następujących skrótów:
\begin{enumerate}[label=(\roman*)]
\item $\forall x\,\psi(x)$ oznacza: „dla każdego zbioru $x$
      zachodzi $\psi(x)$”; $\exists x\,\psi(x)$: „istnieje zbiór $x$”.
\item $\forall x\in a\;\psi(x)$ jest skrótem formuły
      $\forall x\,(x\in a\to\psi(x))$.
      Analogicznie $\exists x\in a\;\psi(x)$ oznacza
      $\exists x\,(x\in a\land\psi(x))$.
\item $x\subseteq a$ oznacza
      $\forall t\,(t\in x\to t\in a)$; nie jest osobnym
      symbolem pierwotnego języka ZF.
\item $\exists!y\,\psi(y)$ oznacza „istnieje dokładnie
      jeden $y$ spełniający $\psi$”, czyli
      $\exists y\,[\psi(y)\land\forall z\,(\psi(z)\to z=y)]$.
\item $\vec p=(p_1,\ldots,p_n)$ oznacza skończony ciąg
      \emph{parametrów}, które występują w rozważanej formule;
      $\varphi(x,\vec p)$ nie jest dodatkowym spójnikiem,
      lecz sposobem pokazania, od jakich zmiennych zależy formuła.
\item $\{x\in a:\varphi(x)\}$, $\bigcup a$, $\mathcal P(a)$
      oraz $\varnothing$ są nazwami \emph{skonstruowanych zbiorów}.
      W pierwotnym języku mającym jedynie $\in$ i równość
      trzeba je rozwinąć do formuł określających ich elementy.
\end{enumerate}
\end{definicja}

\begin{przyklad}[Czytanie kwantyfikatorów od lewej do prawej]
Rozważmy formułę
\[
\forall a\,\forall b\,\exists p\,\forall x\;
\bigl(x\in p\leftrightarrow(x=a\lor x=b)\bigr).
\]
Przeczytajmy każdy fragment. „Dla \emph{dowolnego} zbioru $a$
i dla \emph{dowolnego} zbioru $b$ istnieje zbiór $p$,
taki że dla \emph{każdego} zbioru $x$ zachodzi:
$x$ jest elementem $p$ dokładnie wtedy, gdy $x=a$ lub $x=b$”.
W szczególności $p=\{a,b\}$. Warto zwrócić uwagę, że
$\forall x$ jest objęte zakresem $\exists p$: po ustaleniu
$a,b$ wybieramy \emph{jeden} zbiór $p$ działający dla
wszystkich $x$, a nie osobne $p$ dla każdego $x$.
\end{przyklad}

\begin{przyklad}[Zamiana liter nie zmienia aksjomatu]
Zapis $\forall x\,\forall y\,\forall z\;\Theta(x,y,z)$
i zapis $\forall a\,\forall b\,\forall t\;\Theta(a,b,t)$
wyrażają to samo zdanie, jeżeli w drugim zapisie
w całej formule jednocześnie i konsekwentnie zmienimy
nazwy \emph{związanych} zmiennych. To jest zmiana nazw
zmiennych związanych, a nie zmiana treści twierdzenia.
Dlatego w części formalnej zwykle stosujemy małe litery,
a wielkie litery pozostawiamy na intuicyjne nazwy
ustalonych zbiorów w przykładach.
\end{przyklad}

\begin{uwaga}[Zbiory w teorii a uniwersum modelu]
W modelu $\mathfrak M=(M,E)$ teorii mnogości relacja $E$ interpretuje
symbol $\in$. Elementy $M$ są obiektami nazywanymi „zbiorami”
\emph{wewnątrz modelu}; nie należy automatycznie utożsamiać $E$
z zewnętrzną relacją przynależności używaną w metateorii.
Tę różnicę zilustrowaliśmy w Twierdzeniu~\ref{thm:skolem-paradox}.
\end{uwaga}

\section{Ekstensjonalność, pustka i pary}
\begin{definicja}[Aksjomat ekstensjonalności]\label{def:zf-extensionality}
Postulujemy zdanie
\[
\forall x\,\forall y\Bigl(\bigl(\forall z\,(z\in x\leftrightarrow z\in y)\bigr)
\to x=y\Bigr).
\]
Dwa zbiory o tych samych elementach są więc równe.
\end{definicja}

\begin{definicja}[Aksjomat istnienia zbioru pustego]\label{def:zf-empty}
Postulujemy $\exists e\,\forall x\,(x\notin e)$.
Taki zbiór jest jedyny na mocy ekstensjonalności;
oznaczamy go $\varnothing$.
\end{definicja}

\begin{definicja}[Aksjomat pary]\label{def:zf-pairing}
Dla dowolnych zbiorów $a,b$ istnieje zbiór $p$ taki, że
\[
\forall x\bigl(x\in p\leftrightarrow(x=a\lor x=b)\bigr).
\]
Jedyny taki zbiór zapisujemy $\{a,b\}$. W szczególności
$\{a\}:=\{a,a\}$ oznacza singleton.
\end{definicja}


\begin{przyklad}[Rozwinięcie aksjomatu ekstensjonalności]
Aksjomat~\ref{def:zf-extensionality} ma budowę:
\[
\forall x\,\forall y\;
\underbrace{\Bigl(\forall z\,(z\in x\leftrightarrow z\in y)\Bigr)}_{\text{te same elementy}}
\to\underbrace{(x=y)}_{\text{równe zbiory}}.
\]
Najpierw ustalamy dwa dowolne zbiory $x$ i $y$.
Następnie rozpatrujemy \emph{warunek}: każdy obiekt $z$
należy do $x$ dokładnie wtedy, gdy należy do $y$.
Dopiero z tego warunku wnioskujemy, że $x=y$.
Implikacja nie głosi, że dowolne dwa zbiory są równe.
Przykładowo $x=\{\varnothing\}$ i $y=\varnothing$
nie spełniają poprzednika, bo $\varnothing\in x$, lecz
$\varnothing\notin y$.
\end{przyklad}

\begin{przyklad}[Nieuporządkowana para a para uporządkowana]
Zbiór $\{a,b\}$ nie rozróżnia kolejności, bo $\{a,b\}=\{b,a\}$.
Natomiast para Kuratowskiego
$(a,b):=\{\{a\},\{a,b\}\}$ rozróżnia współrzędne:
jak wykazaliśmy w rozdziale~\ref{ch:relations},
$(a,b)=(c,d)$ implikuje $a=c$ i $b=d$.
Aksjomat pary zapewnia, że oba zbiory występujące w tym zapisie
rzeczywiście istnieją.
\end{przyklad}


\begin{przyklad}[Aksjomat pary, krok po kroku]
Ustalmy $a=\varnothing$ i $b=\{\varnothing\}$.
Aksjomat pary gwarantuje istnienie $p$ spełniającego
\[
\forall x\quad(x\in p\leftrightarrow
  (x=\varnothing\lor x=\{\varnothing\})).
\]
Sprawdźmy, co to oznacza:
\begin{enumerate}[label=(\roman*)]
\item biorąc $x=\varnothing$, widzimy, że $x\in p$;
\item biorąc $x=\{\varnothing\}$, także otrzymujemy $x\in p$;
\item jeżeli $x$ nie jest żadnym z tych dwóch zbiorów,
      prawa strona równoważności jest fałszywa,
      więc $x\notin p$.
\end{enumerate}
Zatem $p=\{\varnothing,\{\varnothing\}\}$.
Aksjomat nie wymaga, aby $a\ne b$; dla $a=b$
daje zbiór jednoelementowy $\{a\}$.
\end{przyklad}

\section{Suma zbioru, zbiór potęgowy i separacja}
\begin{definicja}[Aksjomat sumy]\label{def:zf-union}
Dla każdego zbioru $a$ istnieje zbiór $u$ spełniający
\[
\forall x\bigl(x\in u\leftrightarrow
\exists b\,(b\in a\land x\in b)\bigr).
\]
Piszemy $u=\bigcup a$. Dla $a=\{A,B\}$ otrzymujemy
zwykłą sumę dwóch zbiorów $A\cup B$.
\end{definicja}

\begin{definicja}[Aksjomat zbioru potęgowego]\label{def:zf-power}
Dla każdego zbioru $a$ istnieje zbiór $p$ taki, że
\[
\forall x\bigl(x\in p\leftrightarrow x\subseteq a\bigr),\qquad
x\subseteq a\ \Longleftrightarrow\ \forall t\,(t\in x\to t\in a).
\]
Oznaczamy go $\mathcal P(a)$; zawiera wszystkie \emph{zbiory}
będące podzbiorami $a$.
\end{definicja}


\begin{przyklad}[Aksjomat sumy na konkretnym zbiorze]
Niech $a=\{\{\varnothing\},\{\{\varnothing\}\}\}$.
Jego elementami są dwa zbiory:
$b_1=\{\varnothing\}$ i $b_2=\{\{\varnothing\}\}$.
Aksjomat sumy mówi, że
\[
\bigcup a=\{x:\exists b\in a\ (x\in b)\}.
\]
Jedyny element $b_1$ to $\varnothing$, a jedyny
element $b_2$ to $\{\varnothing\}$. Zatem
$\bigcup a=\{\varnothing,\{\varnothing\}\}$.
Nie chodzi o sumowanie liczb: symbol $\bigcup$ oznacza
zebranie \emph{wszystkich elementów zbiorów należących do $a$}.
\end{przyklad}

\begin{definicja}[Schemat wyróżniania (separacji)]\label{def:zf-separation}
Niech $\varphi(x,\vec p)$ będzie dowolną formułą języka teorii
mnogości, której parametry $\vec p$ są ustalonymi zbiorami.
Dla każdego zbioru $a$ istnieje zbiór $b$ o własności
\[
\forall x\bigl(x\in b\leftrightarrow
(x\in a\land\varphi(x,\vec p))\bigr).
\]
Oznaczamy go $\{x\in a:\varphi(x,\vec p)\}$.
Jedno \emph{wystąpienie schematu} powstaje po ustaleniu jednej
formuły $\varphi$; mówimy o schemacie, a nie pojedynczym
zdaniu, ponieważ takich formuł jest nieskończenie wiele.
\end{definicja}

\begin{przyklad}[Dlaczego paradoks Russella już nie powstaje?]
Dla dowolnego \emph{wcześniej danego} zbioru $a$ wolno zbudować
$R_a:=\{x\in a:x\notin x\}$. Gdyby $R_a\in a$,
to z definicji otrzymalibyśmy
$R_a\in R_a\iff R_a\notin R_a$, sprzeczność.
Wniosek jest więc inny: $R_a\notin a$. Tym samym
nie może istnieć zbiór uniwersalny $V$ zawierający wszystkie
zbiory, bo musiałby zawierać $R_V$. Separacja nie pozwala
wskazać „zbioru wszystkich”, zanim nie dostaliśmy zbioru
ograniczającego $a$.
\end{przyklad}

\begin{twierdzenie}[Przecięcie i różnica są zbiorami]\label{thm:zf-intersections}
Dla każdych zbiorów $A,B$ istnieją zbiory $A\cap B$ i
$A\setminus B$. Dla każdej niepustej rodziny $\mathcal F$
będącej zbiorem istnieje $\bigcap\mathcal F$.
\end{twierdzenie}
\begin{proof}
Z separacji dla formuły $x\in B$ otrzymujemy
$A\cap B=\{x\in A:x\in B\}$. Dla formuły $x\notin B$
otrzymujemy $A\setminus B=\{x\in A:x\notin B\}$.
Jeżeli $\mathcal F\ne\varnothing$, wybieramy jeden element
$F_0\in\mathcal F$. Wówczas zbiór
$\{x\in F_0:\forall F\in\mathcal F\ (x\in F)\}$
istnieje dzięki separacji i jest równy przecięciu rodziny.
Nie potrzebujemy aksjomatu wyboru: wybieramy jeden element
z jednego niepustego zbioru.
\end{proof}

\begin{twierdzenie}[Iloczyn kartezjański jest zbiorem]\label{thm:zf-product}
Dla dowolnych zbiorów $A,B$ istnieje zbiór
$A\times B=\{(a,b):a\in A,\ b\in B\}$.
\end{twierdzenie}
\begin{proof}
Dla $(a,b)=\{\{a\},\{a,b\}\}$ i $a\in A$, $b\in B$
wszystkie elementy tej pary należą do $\mathcal P(A\cup B)$.
Stąd $(a,b)\in\mathcal P(\mathcal P(A\cup B))$.
Korzystając dwukrotnie z aksjomatu zbioru potęgowego,
a następnie z separacji dla formuły
$\exists a\in A\,\exists b\in B\ (z=(a,b))$,
wyodrębniamy wszystkie i tylko takie pary.
\end{proof}


\begin{przyklad}[Schemat wyróżniania jest ograniczony zbiorem]
Niech $a=\{\varnothing,\{\varnothing\},\{\{\varnothing\}\}\}$
i niech $\varphi(x)$ oznacza formułę „$\varnothing\in x$”.
Wówczas separacja tworzy zbiór
\[
 b=\{x\in a:\varnothing\in x\}=\{\{\varnothing\}\}.
\]
Sprawdźmy kolejno trzy elementy $a$:
$\varnothing$ nie ma elementów, więc odpada;
$\{\varnothing\}$ zawiera $\varnothing$, więc zostaje;
$\{\{\varnothing\}\}$ zawiera wyłącznie $\{\varnothing\}$,
więc również odpada. Widzimy, że \emph{najpierw istnieje}
zbiór $a$, a dopiero potem wolno wyodrębnić jego elementy.
\end{przyklad}

\section{Zastępowanie i nieskończoność}
\begin{definicja}[Schemat zastępowania]\label{def:zf-replacement}
Dla każdej formuły $\varphi(x,y,\vec p)$ zachodzi następujący
aksjomat: jeśli dla każdego $x\in a$ istnieje dokładnie jeden
$y$ taki, że $\varphi(x,y,\vec p)$, to istnieje zbiór
\[
 b=\{y:\exists x\in a\ \varphi(x,y,\vec p)\}.
\]
Ściśle: dla każdej formuły $\varphi$ schemat zawiera domknięcie ogólne
implikacji $[\forall x\in a\;\exists!y\,\varphi(x,y,\vec p)]\Rightarrow
[\exists b\;\forall y\,(y\in b\leftrightarrow\exists x\in a\;\varphi(x,y,\vec p))]$.
Formuła może używać parametrów, ale musi wyznaczać
\emph{jednoznaczną} wartość na każdym elemencie $a$.
\end{definicja}

\begin{przyklad}[Obraz zbioru przez funkcję]
Jeśli $f:A\to B$ jest funkcją i $C\subseteq A$, to
$f[C]=\{f(x):x\in C\}$ istnieje dzięki zastępowaniu,
bo przyporządkowanie $x\mapsto f(x)$ jest jednoznaczne.
W tym prostym przypadku można też użyć separacji w $B$:
$f[C]=\{y\in B:\exists x\in C\ (f(x)=y)\}$.
Schemat zastępowania jest szczególnie użyteczny wtedy,
gdy nie znamy z góry zbioru obejmującego wszystkie wyniki.
\end{przyklad}

\begin{definicja}[Aksjomat nieskończoności]\label{def:zf-infinity}
Istnieje zbiór $I$ taki, że $\varnothing\in I$ i dla
każdego $x\in I$ zachodzi $x\cup\{x\}\in I$.
Zapis $S(x):=x\cup\{x\}$ nazywamy tu \emph{następnikiem zbioru}.
Zbiór spełniający te dwa warunki nazywamy \emph{indukcyjnym}.
\end{definicja}

\begin{twierdzenie}[Istnienie najmniejszego zbioru indukcyjnego]
\label{thm:zf-omega}
W ZF istnieje zbiór $\omega$ zawierający $\varnothing$
i zamknięty na $S$, zawarty w każdym zbiorze indukcyjnym.
Ponadto zachodzi zasada indukcji: jeśli $X\subseteq\omega$
zawiera $\varnothing$ i jest zamknięty na $S$, to $X=\omega$.
\end{twierdzenie}
\begin{proof}
Niech $I$ będzie zbiorem indukcyjnym danym aksjomatem
nieskończoności. Zbiór wszystkich indukcyjnych podzbiorów
$I$ istnieje przez separację w $\mathcal P(I)$:
\[
 \mathcal D=\{J\subseteq I:\varnothing\in J\land
 \forall x\in J\ (S(x)\in J)\}.
\]
Rodzina $\mathcal D$ nie jest pusta, bo $I\in\mathcal D$.
Przez Twierdzenie~\ref{thm:zf-intersections} istnieje
$\omega:=\bigcap\mathcal D$. Ponieważ wszystkie elementy
$\mathcal D$ zawierają $\varnothing$ i są zamknięte na $S$,
również ich przecięcie ma obie własności. Jeżeli $J$ jest
dowolnym zbiorem indukcyjnym, to $J\cap I\in\mathcal D$,
a więc $\omega\subseteq J\cap I\subseteq J$.
Wreszcie każdy indukcyjny podzbiór $X$ zbioru $\omega$
jest również indukcyjnym zbiorem, więc z minimalności
$\omega\subseteq X$. Ponieważ $X\subseteq\omega$, mamy równość.
\end{proof}

\begin{uwaga}[Liczby naturalne jako zbiory]
Definiując $0:=\varnothing$, $1:=S(0)=\{0\}$,
$2:=S(1)=\{0,1\}$ itd., otrzymujemy standardową
konstrukcję von Neumanna. W tym sensie zwykła $\N$
z poprzednich rozdziałów może być reprezentowana przez $\omega$.
Nie jest to nowy aksjomat, lecz konstrukcja oparta na ZF.
\end{uwaga}

\section{Fundowanie i twierdzenie Cantora}
\begin{definicja}[Aksjomat ufundowania (regularności)]
\label{def:zf-foundation}
Każdy niepusty zbiór $a$ ma element $b$, który nie ma
żadnego wspólnego elementu z $a$:
\[
 \forall a\bigl(a\ne\varnothing\to
 \exists b\in a\ (a\cap b=\varnothing)\bigr).
\]
Jest to aksjomat ufundowania. W szczególności wyklucza
$x\in x$: w przeciwnym razie zbiór $\{x\}$
nie miałby elementu rozłącznego z $\{x\}$.
\end{definicja}

\begin{twierdzenie}[Twierdzenie Cantora]\label{thm:zf-cantor}
Dla każdego zbioru $A$ nie istnieje surjekcja
$f:A\to\mathcal P(A)$. Jednocześnie istnieje iniekcja
$\iota:A\to\mathcal P(A)$ zadana przez $\iota(a)=\{a\}$.
\end{twierdzenie}
\begin{proof}
Przypuśćmy, że $f:A\to\mathcal P(A)$ jest surjekcją.
Z separacji istnieje zbiór
\[
 D=\{a\in A:a\notin f(a)\}\subseteq A.
\]
Ponieważ $D\in\mathcal P(A)$, z surjektywności istnieje
$d\in A$ z $f(d)=D$. Wtedy
$ d\in D\iff d\notin f(d)=D$, sprzeczność.
Dla iniekcji zauważmy, że jeśli $\{a\}=\{b\}$,
to z definicji singletonu $a=b$. Aksjomat pary
zapewnia istnienie singletonów. Zastępowanie zastosowane do
jednoznacznej formuły $y=\{x\}$ gwarantuje istnienie
zbioru $\{\{a\}:a\in A\}$, będącego obrazem tej iniekcji.
Dowód obejmuje także $A=\varnothing$: surjekcja
$\varnothing\to\mathcal P(\varnothing)=\{\varnothing\}$
nie istnieje, zaś pusta funkcja jest iniekcją.
\end{proof}


\section{Liczby naturalne von Neumanna: konstrukcja i aksjomaty Peana}

Samo istnienie zbioru $\omega$ z Twierdzenia~\ref{thm:zf-omega}
nie daje jeszcze definicji arytmetycznych działań.
Skonstruujemy je w ZF, używając indukcji i zastępowania,
bez aksjomatu wyboru.

\begin{definicja}[Liczby naturalne von Neumanna]
\label{def:zf-vn-natural}
Liczbami naturalnymi nazywamy elementy $\omega$.
Przyjmujemy $0:=\varnothing$, $1:=S(0)$, $2:=S(1)$
oraz $3:=S(2)$, gdzie $S(n):=n\cup\{n\}$. Stąd
$0=\varnothing$, $1=\{0\}$, $2=\{0,1\}$ i $3=\{0,1,2\}$.
\end{definicja}

\begin{lemat}[Własności następnika]\label{lem:zf-vn-successor}
Dla wszystkich $m,n\in\omega$:
\begin{enumerate}[label=(\roman*)]
\item zbiór $n$ jest przechodni, tj. $x\in y\in n$
      implikuje $x\in n$; ponadto $n\subseteq\omega$ i $n\notin n$;
\item $S(n)\ne0$, zaś $S(m)=S(n)$ implikuje $m=n$;
\item każdy $n\in\omega$ jest zerem albo ma postać $S(k)$,
      gdzie $k\in\omega$;
\item zbiór $\omega$ jest przechodni.
\end{enumerate}
\end{lemat}
\begin{proof}
Przechodniość $n$ dowodzimy indukcją po $n$.
Zbiór $0$ jest przechodni. Jeśli $n$ jest przechodni,
a $x\in y\in S(n)$, to albo $y=n$,
albo $y\in n$. W pierwszym przypadku $x\in n\subseteq S(n)$,
a w drugim z przechodniości $n$ także $x\in n\subseteq S(n)$.
Ufundowanie daje $n\notin n$: w przeciwnym razie
jedyny element zbioru $\{n\}$ przecinałby $\{n\}$.

Ponieważ $n\in S(n)$, zbiór $S(n)$ jest niepusty.
Jeśli $S(m)=S(n)$ i $m\ne n$, to
$m\in S(n)$ i $n\in S(m)$ prowadzą do
$m\in n$ oraz $n\in m$.
Żaden z elementów zbioru $\{m,n\}$ nie byłby
z nim rozłączny, wbrew aksjomatowi ufundowania.

Dla (iii) schemat zastępowania daje zbiór $S[\omega]$.
Zbiór $J:=\{0\}\cup S[\omega]$ zawiera $0$
i jest zamknięty na $S$: jeżeli $k\in J\subseteq\omega$,
to $S(k)\in S[\omega]\subseteq J$.
Z minimalności $\omega$ wynika $J=\omega$.

Wreszcie przez indukcję wykażemy $n\subseteq\omega$
dla wszystkich $n\in\omega$. Dla $n=0$ jest to jasne;
jeśli $n\subseteq\omega$, to
$S(n)=n\cup\{n\}\subseteq\omega$, ponieważ $n\in\omega$.
Jest to również przechodniość całego $\omega$.

Sprawdźmy jeszcze własność używaną w konstrukcji rekurencyjnej:
$0\in S(n)$ dla każdego $n\in\omega$.
Dla $n=0$ mamy $S(0)=\{0\}$.
Jeśli $0\in S(n)$, to
$S(n)\subseteq S(S(n))$, zatem $0\in S(S(n))$.
Indukcja na $\omega$ kończy uzasadnienie.
Ponadto $n\in S(n)$ bezpośrednio z definicji następnika.
\end{proof}

\begin{wniosek}[Aksjomaty Peana dla $(\omega,0,S)$]
\label{cor:zf-peano}
W strukturze $(\omega,0,S)$ zachodzą:
$0\in\omega$, $S[\omega]\subseteq\omega$,
$0\notin S[\omega]$, $S$ jest iniekcją,
a każdy podzbiór $X\subseteq\omega$ zawierający
$0$ i zamknięty na $S$ jest równy $\omega$.
Podzbiór oznacza tu zbiór w sensie ZF;
nie twierdzimy, że w dowolnym modelu ZF
indukcja obejmuje wszystkie zewnętrzne podzbiory
jego interpretacji $\omega$.
\end{wniosek}
\begin{proof}
Pierwsze dwa warunki daje konstrukcja
zbioru indukcyjnego. Dwa następne
wynikają z Lematu~\ref{lem:zf-vn-successor}(ii),
a ostatni z Twierdzenia~\ref{thm:zf-omega}.
Nie wymaga to kwantyfikacji po wszystkich
zewnętrznych podzbiorach uniwersum modelu.
\end{proof}

\begin{przyklad}[Pierwsze następniki]
\label{ex:zf-vn-steps}
Z definicji mamy $S(2)=2\cup\{2\}=\{0,1,2\}=3$.
Zachodzi $2\in3$ i $2\subseteq3$,
natomiast $3\notin3$ na mocy ufundowania.
\end{przyklad}

\begin{figure}[htbp]
\centering
$\displaystyle
0=\varnothing\ \xrightarrow{\ S\ }\ 
1=\{0\}\ \xrightarrow{\ S\ }\ 
2=\{0,1\}\ \xrightarrow{\ S\ }\ 
3=\{0,1,2\}.
$
\caption{Pierwsze liczby naturalne von Neumanna i operacja następnika.}
\label{fig:zf-vn-successor}
\end{figure}

\section{Twierdzenie o rekursji na $\omega$ w ZF}

\begin{twierdzenie}[Rekursja na liczbach naturalnych]
\label{thm:zf-omega-recursion}
Niech $B$ będzie niepustym zbiorem,
$b\in B$ oraz $G:\omega\times B\to B$ funkcją.
Istnieje dokładnie jedna funkcja $F:\omega\to B$
taka, że dla każdego $n\in\omega$
\[
F(0)=b,\qquad F(S(n))=G(n,F(n)).
\]
Twierdzenie jest dowodliwe w ZF, bez AC.
\end{twierdzenie}
\begin{proof}
Dla $n\in\omega$ nazwijmy \emph{$n$-przybliżeniem}
funkcję $f:S(n)\to B$ spełniającą $f(0)=b$
oraz, dla każdego $k\in S(n)$, warunek
$S(k)\in S(n)\Longrightarrow
f(S(k))=G(k,f(k))$.
Warunek $f(0)=b$ ma sens, ponieważ
$0\in S(n)$ dla każdego $n\in\omega$,
co wykazaliśmy w dowodzie
Lematu~\ref{lem:zf-vn-successor}.
Ponieważ $S(n)\times B$ jest zbiorem,
wszystkie takie funkcje tworzą zbiór
przez separację w $\mathcal P(S(n)\times B)$.
W szczególności nie zakładamy wcześniej
istnienia funkcji określonej na całym $\omega$.

Indukcyjnie po $n$ dowodzimy, że istnieje dokładnie
jedno $n$-przybliżenie $f_n$.
Dla $n=0$ jest nim $\{(0,b)\}$.
Jeśli istnieje dokładnie jedno $f_n$, dołączamy
parę $(S(n),G(n,f_n(n)))$.
Wartość $f_n(n)$ jest określona,
gdyż $n\in S(n)$ zgodnie z definicją następnika.
Nowy argument nie należał do $S(n)$,
gdyż $S(n)\notin S(n)$,
a $S(S(n))=S(n)\cup\{S(n)\}$.
Powstaje $(S(n))$-przybliżenie $f_{S(n)}$.
Dowolne inne przybliżenie na $S(S(n))$
po ograniczeniu do $S(n)$ musi równać się $f_n$;
jego jedyna nowa wartość jest wyznaczona równaniem
rekurencyjnym, więc również jest jedyna.

Jednoznacznie określona wartość $f_n(n)$
jest definiowalna formułą teorii mnogości
z parametrami $B,b,G$: piszemy, że istnieje
funkcja $f$ o dziedzinie $S(n)$ spełniająca
warunek przybliżenia oraz $f(n)=y$.
Poprzednia indukcja gwarantuje dla każdego
$n\in\omega$ istnienie dokładnie jednego takiego $y$.
Schemat zastępowania, zastosowany do jednoznacznego
przyporządkowania $n\mapsto (n,y)$,
daje zbiór par
$F:=\{(n,f_n(n)):n\in\omega\}$.
To funkcja $\omega\to B$, przy czym
$F(0)=b$ oraz
$F(S(n))=f_{S(n)}(S(n))=G(n,f_n(n))=G(n,F(n))$.
Jeżeli $F,H$ spełniają oba równania,
to zbiór $\{n\in\omega:F(n)=H(n)\}$
istnieje przez separację, zawiera $0$
i jest zamknięty na $S$. Indukcja w $\omega$
daje $F=H$.
\end{proof}

\begin{uwaga}[Rekursja nie wymaga wyboru]
\label{rem:zf-recursion-not-choice}
Każda wartość przybliżenia jest jednoznacznie
wyznaczona, a schemat zastępowania tworzy
jeden zbiór wszystkich par. Nie wybieramy
dowolnie elementów z rodziny niepustych zbiorów.
\end{uwaga}

\section{Dodawanie: definicja i wyprowadzenie praw}

\begin{definicja}[Dodawanie w $\omega$]\label{def:zf-addition}
Dla każdego $a\in\omega$ zastosujmy
Twierdzenie~\ref{thm:zf-omega-recursion}
z $B=\omega$, wartością początkową $a$
i funkcją $G(n,z)=S(z)$.
Definiujemy $a+n$ przez
\begin{equation}\label{eq:zf-add}
a+0=a,\qquad a+S(n)=S(a+n).
\end{equation}
Schemat zastępowania zapewnia,
że wykres funkcji $+:\omega\times\omega\to\omega$
jest zbiorem.
\end{definicja}

\begin{lemat}[Zero i następnik w dodawaniu]
\label{lem:zf-add-basic}
Dla wszystkich $a,n\in\omega$:
$0+n=n$ oraz $S(a)+n=S(a+n)$.
\end{lemat}
\begin{proof}
Pierwszą tożsamość dowodzimy indukcją po $n$.
Dla $n=0$ jest jasna; z $0+n=n$
wynika $0+S(n)=S(0+n)=S(n)$.
Druga tożsamość dla $n=0$ ma postać
$S(a)=S(a+0)$.
Jeśli $S(a)+n=S(a+n)$, to
\[
S(a)+S(n)=S(S(a)+n)=S(S(a+n))=S(a+S(n)).
\]
Każde przejście wynika z definicji dodawania
albo z hipotezy indukcyjnej.
\end{proof}

\begin{twierdzenie}[Łączność i przemienność dodawania]
\label{thm:zf-add-assoc-comm}
Dla wszystkich $a,b,c\in\omega$:
$(a+b)+c=a+(b+c)$ oraz $a+b=b+a$.
\end{twierdzenie}
\begin{proof}
Ustalmy $a,b$ i indukujmy po $c$.
Dla $c=0$ obie strony pierwszej równości
to $a+b$. W kroku indukcyjnym:
\[
(a+b)+S(c)=S((a+b)+c)=S(a+(b+c))
=a+S(b+c)=a+(b+S(c)).
\]
Przemienność dowodzimy indukcją po $b$ przy
ustalonym $a$. Dla $b=0$ korzystamy z
$a+0=a=0+a$. Jeśli $a+b=b+a$, to
\[
a+S(b)=S(a+b)=S(b+a)=S(b)+a,
\]
gdzie ostatnia równość pochodzi
z Lematu~\ref{lem:zf-add-basic}.
\end{proof}

\begin{lemat}[Suma zerowa i skracanie]
\label{lem:zf-add-cancel}
Dla $a,b,c\in\omega$:
\begin{enumerate}[label=(\roman*)]
\item $a+b=0$ wtedy i tylko wtedy, gdy $a=b=0$;
\item $a+c=b+c$ implikuje $a=b$.
\end{enumerate}
\end{lemat}
\begin{proof}
Jeśli $b\ne0$, to z Lematu~\ref{lem:zf-vn-successor}(iii)
mamy $b=S(k)$, a wtedy
$a+b=S(a+k)\ne0$. Gdy $b=0$,
równość $a+b=0$ oznacza $a=0$.
Kierunek odwrotny jest oczywisty.
Dla skracania indukujemy po $c$.
Przy $c=0$ teza jest bezpośrednia.
Jeśli $a+S(c)=b+S(c)$, to
$S(a+c)=S(b+c)$, więc
iniekcyjność $S$ daje $a+c=b+c$,
a hipoteza indukcyjna --- $a=b$.
\end{proof}

\section{Mnożenie, rozdzielność i potęgowanie}

\begin{definicja}[Mnożenie w $\omega$]\label{def:zf-multiplication}
Dla każdego $a\in\omega$ z
Twierdzenia~\ref{thm:zf-omega-recursion},
przy $B=\omega$, $b=0$ i $G(n,z)=z+a$,
otrzymujemy jednoznaczną funkcję spełniającą
\begin{equation}\label{eq:zf-mult}
a\cdot0=0,\qquad a\cdot S(n)=a\cdot n+a.
\end{equation}
Jej dwuzmienny wykres istnieje przez zastępowanie.
\end{definicja}

\begin{lemat}[Zero i jedynka w mnożeniu]
\label{lem:zf-mult-basic}
Dla każdego $a\in\omega$ mamy
$a\cdot0=0$, $0\cdot a=0$,
$a\cdot1=a$ i $1\cdot a=a$.
\end{lemat}
\begin{proof}
Pierwsza równość jest definicją,
drugą uzyskujemy przez indukcję po $a$:
$0\cdot S(a)=0\cdot a+0=0$.
Dalej
$a\cdot1=a\cdot S(0)=0+a=a$.
Dla ostatniej równości indukujemy po $a$:
$1\cdot0=0$, a z $1\cdot a=a$ dostajemy
$1\cdot S(a)=1\cdot a+1=a+S(0)=S(a)$.
\end{proof}

\begin{twierdzenie}[Rozdzielność mnożenia względem dodawania]
\label{thm:zf-distributivity}
Dla wszystkich $a,b,c\in\omega$:
$a(b+c)=ab+ac$ i $(a+b)c=ac+bc$.
\end{twierdzenie}
\begin{proof}
W obu przypadkach indukujemy po $c$.
Dla $c=0$ stosujemy równości zerowe.
Jeśli $a(b+c)=ab+ac$, to
\[
a(b+S(c))=aS(b+c)=a(b+c)+a
=(ab+ac)+a=ab+(ac+a)=ab+aS(c).
\]
Druga tożsamość w kroku indukcyjnym
wynika z równości
\[
(a+b)S(c)=(a+b)c+(a+b)
=(ac+bc)+(a+b)=(ac+a)+(bc+b)=aS(c)+bS(c).
\]
W przestawianiu składników używamy
Twierdzenia~\ref{thm:zf-add-assoc-comm}.
\end{proof}

\begin{twierdzenie}[Przemienność i łączność mnożenia]
\label{thm:zf-mult-assoc-comm}
Dla wszystkich $a,b,c\in\omega$:
$ab=ba$ oraz $(ab)c=a(bc)$.
\end{twierdzenie}
\begin{proof}
Dla przemienności ustalmy $a$ i indukujmy po $b$.
Baza to $a0=0=0a$. W kroku:
\[
aS(b)=ab+a=ba+1a=(b+1)a=S(b)a,
\]
przy czym korzystamy z rozdzielności
i Lematu~\ref{lem:zf-mult-basic}.
Dla łączności ustalmy $a,b$
i indukujmy po $c$. Baza jest zerowa,
a krok ma postać
\[
(ab)S(c)=(ab)c+ab=a(bc)+ab
=a(bc+b)=a(bS(c)).
\]
Użyliśmy wyłącznie już wykazanej rozdzielności.
\end{proof}

\begin{definicja}[Potęgowanie]
\label{def:zf-exponentiation}
Dla $a\in\omega$ rekurencja z
Twierdzenia~\ref{thm:zf-omega-recursion},
przy $B=\omega$, $b=1$ i $G(n,z)=z\cdot a$,
wyznacza $a^n$ przez
$a^0=1$ i $a^{S(n)}=a^n\cdot a$.
W szczególności $0^0=1$ jest konwencją
tej konstrukcji.
\end{definicja}

\begin{twierdzenie}[Prawo wykładników]
\label{thm:zf-exponent-add}
Dla każdego $a,m,n\in\omega$
zachodzi $a^{m+n}=a^m\cdot a^n$.
\end{twierdzenie}
\begin{proof}
Indukujemy po $n$. Dla $n=0$
obie strony wynoszą $a^m$,
gdyż $a^0=1$.
Z hipotezy indukcyjnej:
\[
a^{m+S(n)}=a^{S(m+n)}=a^{m+n}\cdot a
=(a^m a^n)a=a^m(a^n a)=a^m a^{S(n)}.
\]
Łączność mnożenia pochodzi z
Twierdzenia~\ref{thm:zf-mult-assoc-comm}.
\end{proof}

\begin{przyklad}[Obliczanie z definicji]
\label{ex:zf-arithmetic}
Mamy $2+1=2+S(0)=S(2)=3$,
a $2\cdot2=2\cdot S(1)=2\cdot1+2=2+2=4$.
Kolejno $2^0=1$, $2^1=2$, $2^2=4$, $2^3=8$.
Ostatnie symbole oznaczają odpowiednie
iteracje następnika, nie nowe obiekty pierwotne.
\end{przyklad}

\section{Porządek naturalny, minimum i indukcja silna}

\begin{definicja}[Porządek arytmetyczny]
\label{def:zf-natural-order}
Dla $m,n\in\omega$ definiujemy
$m\le n$ wtedy i tylko wtedy, gdy
istnieje $k\in\omega$ z $m+k=n$.
Piszemy $m<n$, gdy $m\le n$ i $m\ne n$.
Relacji tej nie utożsamiamy z $\in$
bez osobnego dowodu.
\end{definicja}

\begin{lemat}[Odcinki początkowe $\omega$]
\label{lem:zf-initial-segments}
Dla każdego $n\in\omega$ zachodzi
$\{m\in\omega:m\le n\}=S(n)=n\cup\{n\}$.
W szczególności $m<n$ wtedy i tylko wtedy,
gdy $m\in n$.
\end{lemat}
\begin{proof}
Indukujemy po $n$. Przy $n=0$ równość
$m+k=0$ wymusza $m=k=0$
(Lemat~\ref{lem:zf-add-cancel}),
zatem jedynym $m\le0$ jest $0$.
Jeśli teza zachodzi dla $n$,
a $m+k=S(n)$, to przy $k=0$
mamy $m=S(n)$. Przy $k\ne0$
zapisujemy $k=S(t)$, zatem
$S(m+t)=S(n)$. Iniekcyjność
następnika daje $m+t=n$,
czyli $m\le n$ i na mocy
hipotezy indukcyjnej $m\in S(n)$.
W obu sytuacjach $m\in S(S(n))$.
Odwrotnie, jeśli $m\in S(S(n))$,
to albo $m=S(n)$, i wtedy
$m+0=S(n)$, albo $m\in S(n)$.
W drugim przypadku istnieje $t$
z $m+t=n$, więc $m+S(t)=S(n)$.
To daje $m\le S(n)$.
Ostatnią równoważność otrzymujemy,
usuwając z $S(n)$ jego element $n$.
\end{proof}

\begin{twierdzenie}[Dobry porządek $\omega$]
\label{thm:zf-wellordered-omega}
Relacja $\le$ z Definicji~\ref{def:zf-natural-order}
jest dobrym porządkiem $\omega$:
każdy niepusty podzbiór $\omega$
ma element najmniejszy.
\end{twierdzenie}
\begin{proof}
Zwrotność wynika z $m+0=m$,
a przechodniość z łączności dodawania.
Jeśli $m+k=n$ i $n+t=m$,
to $m+(k+t)=m+0$, zatem
z prawa skracania $k+t=0$.
Stąd $k=t=0$ i $m=n$:
porządek jest antysymetryczny.

Liniowość dowodzimy indukcją po $n$
dla własności $\forall m\in\omega\,
(m\le n\lor n\le m)$.
Przy $n=0$ mamy $0+m=m$.
Jeżeli $m\le n$, to
$m+k=n$ daje $m+S(k)=S(n)$,
zatem $m\le S(n)$. Jeśli $n\le m$,
mamy $m=n+k$. Przy $k=0$
otrzymujemy $m=n\le S(n)$,
a przy $k=S(t)$:
$m=n+S(t)=S(n)+t$, więc $S(n)\le m$.

Niech $A\subseteq\omega$ będzie niepusty.
Z aksjomatu ufundowania wybieramy $m\in A$
z $m\cap A=\varnothing$. Dla każdego $n\in A$
liniowość daje $m\le n$ lub $n\le m$.
Druga możliwość z $n\ne m$
oznacza $n\in m$ na mocy
Lematu~\ref{lem:zf-initial-segments},
sprzecznie z $m\cap A=\varnothing$.
Wobec tego $m\le n$ dla wszystkich $n\in A$.
\end{proof}

\begin{wniosek}[Indukcja silna]\label{cor:zf-strong-induction}
Niech $P(n)$ będzie dowolną formułą ZF
z parametrami i wyróżnioną zmienną $n$.
Jeżeli dla każdego $n\in\omega$
z założenia $\forall k\in n\,P(k)$
wynika $P(n)$, to $P(n)$ zachodzi
dla wszystkich $n\in\omega$.
\end{wniosek}
\begin{proof}
Gdyby zbiór $C:=\{n\in\omega:\neg P(n)\}$
był niepusty, miałby najmniejszy element $n_0$.
Wtedy $P(k)$ dla każdego $k\in n_0$,
co na mocy założenia daje $P(n_0)$,
wbrew $n_0\in C$.
Zbiór $C$ istnieje przez separację.
\end{proof}

\begin{przyklad}[Brak największej liczby naturalnej]
\label{ex:zf-no-largest}
Każde $n\in\omega$ ma następnik
$S(n)\in\omega$, a
$n\in S(n)$ oznacza na mocy
Lematu~\ref{lem:zf-initial-segments},
że $n<S(n)$. Zatem żadne $n$
nie jest największym elementem $\omega$.
\end{przyklad}

\begin{lemat}[Zgodność porządku z działaniami]
\label{lem:zf-natural-monotone}
Dla $a,b,c\in\omega$:
\begin{enumerate}[label=(\roman*)]
\item $a<b$ implikuje $a+c<b+c$;
\item $a<b$ i $c\ne0$ implikują $ac<bc$;
\item $ab=0$ implikuje $a=0$ lub $b=0$;
\item $c\ne0$ i $ac=bc$ implikują $a=b$.
\end{enumerate}
\end{lemat}
\begin{proof}
Jeśli $a<b$, to $b=a+d$ dla pewnego $d\ne0$.
Stąd $b+c=(a+c)+d$, a prawo skracania
wyklucza $b+c=a+c$, co daje (i).
Dla (iii), jeśli $a,b\ne0$,
to $b=S(t)$ i $a=S(k)$,
więc $ab=at+a=at+S(k)=S(at+k)\ne0$.
Dla (ii) z $b=a+d$ mamy
$bc=ac+dc$, gdzie $d\ne0$.
Z (iii) i $c\ne0$ wynika $dc\ne0$,
zatem $ac<bc$.
Dla (iv) liniowość porządku
daje $a=b$, $a<b$ albo $b<a$.
Dwa ostatnie przypadki przeczą (ii)
i równości $ac=bc$.
\end{proof}

\begin{twierdzenie}[Dzielenie z resztą]
\label{thm:zf-division}
Dla wszystkich $a\in\omega$ i
$b\in\omega\setminus\{0\}$
istnieje dokładnie jedna para
$(q,r)\in\omega\times\omega$
taka, że $a=bq+r$ oraz $r<b$.
\end{twierdzenie}
\begin{proof}
Ustalmy $b\ne0$ i zastosujmy indukcję silną
po $a$. Jeśli $a<b$, kładziemy
$q=0$ i $r=a$. Wtedy $a=b\cdot0+a$.
Jeśli $b\le a$, istnieje $t\in\omega$
z $a=b+t$. Ponieważ $0+b=b$,
mamy $0\le b$, a z $b\ne0$ wynika $0<b$.
Z Lematu~\ref{lem:zf-natural-monotone}(i),
po dodaniu $t$, otrzymujemy
$t=t+0<t+b=a$.
Hipoteza indukcyjna daje
$t=bq+r$, gdzie $r<b$,
a wtedy
$a=b+(bq+r)=bS(q)+r$.
Istnienie zostało udowodnione.

Dla jednoznaczności niech
$bq+r=bq'+r'$ i $r,r'<b$.
Jeśli $q=q'$, skracanie dodawania
daje $r=r'$. Przy $q<q'$
możemy zapisać $q'=q+d$
dla $d\ne0$; wtedy $d=S(k)$ i
\[
bq'=bq+bd=bq+(bk+b)\ge bq+b.
\]
Istotnie, $0\le bk$, więc z definicji porządku
$bq+b\le bq+bk+b=bq'$.
Podobnie $0\le r'$, stąd
$bq'\le bq'+r'=a$.
Przechodniość daje $bq+b\le a$.
Z drugiej strony $r<b$ i
Lemat~\ref{lem:zf-natural-monotone}(i)
dają $a=bq+r<bq+b$, sprzeczność.
Przypadek $q'<q$ jest symetryczny.
\end{proof}

\begin{przyklad}[Dzielenie $17$ przez $5$]
\label{ex:zf-division}
Dla liczb von Neumanna $17$ i $5$
mamy $17=5\cdot3+2$ i $2<5$,
ponieważ $2+3=5$.
Twierdzenie~\ref{thm:zf-division}
gwarantuje jednoznaczność $(q,r)=(3,2)$.
\end{przyklad}

\begin{uwaga}[Zakres konstrukcji arytmetyki]
\label{rem:zf-arithmetic-scope}
Dodawanie, mnożenie, potęgowanie
i porządek na $\omega$ zostały skonstruowane
za pomocą aksjomatów ZF, a ich wykresy
są zbiorami. Ani aksjomaty Peana dla
$(\omega,0,S)$, ani twierdzenie o rekursji
nie wymagają aksjomatu wyboru.
Ogólniejszą rekursję pozaskończoną
omówimy przy liczbach porządkowych.
\end{uwaga}

\begin{uwaga}[Czego jeszcze nie przyjęliśmy?]
Żaden z dotychczas wymienionych aksjomatów nie orzeka,
że dla \emph{dowolnej} rodziny niepustych zbiorów można
jednocześnie wybrać po jednym elemencie z każdego z nich.
To będzie osobny aksjomat. Rozdział~\ref{ch:axiom-choice}
wykaże jego równoważność z twierdzeniem o dobrym uporządkowaniu
i z lematem Kuratowskiego--Zorna. Zaznaczmy też, że
mówimy tu o \emph{ZF} (z aksjomatem ufundowania), a nie
o teorii z klasami właściwymi jako osobnymi obiektami.
\end{uwaga}